% 双塔与ANN接口：用户塔请求时算一次 h(x)，物品嵌入离线预计算，内积经近邻索引取Top-K；
% 训练在采样竞争集合上做softmax，与服务面对的全目录不同。跨栏宽图（figure*），内容宽约150mm。
\begin{tikzpicture}[x=1mm,y=1mm,
  ftbox/.style={rectangle,draw=FigureInput!70,fill=FigureInput!10,align=center,
    text=BookInk,rounded corners=1.4pt,minimum height=12mm,text width=26mm,inner sep=2pt},
  ftnet/.style={rectangle,draw=FigureModel!80,thick,fill=FigureModel!12,align=center,
    text=BookInk,rounded corners=1.4pt,minimum height=13mm,text width=26mm,inner sep=2pt},
  ftitem/.style={rectangle,draw=FigureModel!60,fill=FigureModel!7,align=center,
    text=BookInk,rounded corners=1.4pt,minimum height=13mm,text width=26mm,inner sep=2pt},
  ftidx/.style={rectangle,draw=FigureLoss!80,thick,fill=FigureLoss!14,align=center,
    text=BookInk,rounded corners=1.4pt,minimum height=13mm,text width=24mm,inner sep=2pt},
  ftout/.style={rectangle,draw=FigureOutput!75,thick,fill=FigureOutput!14,align=center,
    text=BookInk,rounded corners=1.4pt,minimum height=12mm,text width=20mm,inner sep=2pt},
  ftflow/.style={-{Stealth[length=1.8mm,width=1.2mm]},draw=FigureOutput!85,line width=.9pt},
  ftdash/.style={-{Stealth[length=1.8mm,width=1.2mm]},draw=FigureModel!80,
    line width=.85pt,dashed},
  ftlab/.style={font=\figurefont\scriptsize,text=BookMuted,fill=white,inner sep=1pt}
]
  % 用户塔
  \node[ftbox] (hist) at (18,16)
    {观看、搜索词嵌入汇总\\人口统计等请求特征};
  \node[ftnet] (utower) at (55,16)
    {用户网络\\$\bm h(\bm x)$\\请求时计算一次};
  \draw[ftflow] (hist.east) -- (utower.west);

  % 物品侧
  \node[ftitem] (itable) at (55,-15)
    {物品嵌入表 $\bm v_i$\\离线预计算，可建索引};

  % 内积与ANN
  \node[ftidx] (ann) at (95,16)
    {内积近邻索引\\ANN / 精确扫描};
  \draw[ftflow] (utower.east) -- (ann.west);
  \draw[ftdash] (itable.north) -- (68,-15) -- (68,10) -- (ann.south);

  % 输出
  \node[ftout] (cand) at (132,16) {Top-$K$ 候选};
  \draw[ftflow] (ann.east) -- (cand.west);

  % 训练分支标注
  \node[ftlab,align=left,anchor=west] at (18,-25)
    {训练：下一观看物品为正例，在采样竞争集合上做softmax（见3.4.1.1）；};
  \node[ftlab,align=left,anchor=west] at (18,-29.5)
    {服务：用户向量与全部物品嵌入算内积，训练竞争集合不等于全目录。};
  \node[ftlab,align=center,text=FigureLoss!85] at (95,24)
    {两端只在最后以内积相遇，历史与物品不做深层交互};
\end{tikzpicture}
